% !TEX root = ../main.tex
%% Chapitre 4 — Continuité
%% Sources : frontend/cours/resume/terminal/resume_continuite_fonctions_terminal.txt
%%           frontend/exercices/terminal_spé_Maths/exercice_continuite_fonctions.txt
%% Exercices 6 à 8 : problèmes inédits rédigés pour le livre.
%% Questions avec ln remplacées (ln n'est vu qu'au chapitre 5).

\begin{ouverture}[continuite]{4}{Continuité}{Partie 1 · Analyse}
  {Définition\enspace·\enspace Prolongement\enspace·\enspace Valeurs intermédiaires\enspace·\enspace
   Dichotomie\enspace·\enspace Point fixe}
  \objectif{Étudier la continuité d'une fonction en un point, prolonger une fonction par continuité.}
  \objectif{Appliquer le théorème des valeurs intermédiaires pour justifier l'existence de solutions.}
  \objectif{Démontrer qu'une équation admet une unique solution sur un intervalle et l'encadrer.}
  \objectif{Programmer la méthode de dichotomie.}
  \objectif{Déterminer la limite d'une suite récurrente à l'aide du théorème du point fixe.}
  \prerequis{Limites de fonctions (chapitre 2).}
  \prerequis{Dérivation et sens de variation (chapitre 3).}
  \prerequis{Suites récurrentes et convergence monotone (chapitre 1).}
\end{ouverture}

%% =====================================================================
%%  L'ESSENTIEL
%% =====================================================================
\essentiel

\essSection{1}{Continuité}

\begin{definition}[Fonction continue]
Soit $f$ une fonction définie sur un intervalle $I$ et $a$ un réel de $I$. La fonction $f$ est \textbf{continue en $a$} si
$\displaystyle\lim_{x\to a}f(x)=f(a)$. Elle est \textbf{continue sur $I$} si elle est continue en tout réel de $I$ :
sa courbe se trace alors « sans lever le crayon ».
\end{definition}

\noindent
\begin{minipage}[t]{0.32\linewidth}\centering
\begin{tikzpicture}
\begin{axis}[otSchema, width=54mm, height=40mm, xmin=0, xmax=4, ymin=0, ymax=3,
  xtick={2}, xticklabels={$a$}, ytick=\empty, xlabel={$x$}, xlabel style={at={(1,0)},anchor=west}]
  \addplot[otNavy,line width=1.1pt,domain=0.2:3.8,samples=60] {1.4+0.6*sin(deg(1.4*x))+0.12*x};
  \fill[otNavy] (axis cs:2,1.84) circle (1.5pt);
\end{axis}
\end{tikzpicture}\\[1mm]
\legende{Continue en $a$ :\\ $\lim\limits_{x\to a}f(x)=f(a)$}
\end{minipage}\hfill
\begin{minipage}[t]{0.32\linewidth}\centering
\begin{tikzpicture}
\begin{axis}[otSchema, width=54mm, height=40mm, xmin=0, xmax=4, ymin=0, ymax=3,
  xtick={2}, xticklabels={$a$}, ytick=\empty, xlabel={$x$}, xlabel style={at={(1,0)},anchor=west}]
  \addplot[otNavy,line width=1.1pt,domain=0.2:2] {0.6+0.4*x};
  \addplot[otNavy,line width=1.1pt,domain=2:3.8] {2.2+0.2*(x-2)};
  \draw[otNavy,line width=0.8pt,fill=white] (axis cs:2,1.4) circle (1.8pt);
  \fill[otNavy] (axis cs:2,2.2) circle (1.8pt);
\end{axis}
\end{tikzpicture}\\[1mm]
\legende{Non continue en $a$ :\\ la courbe « saute »}
\end{minipage}\hfill
\begin{minipage}[t]{0.32\linewidth}\centering
\begin{tikzpicture}
\begin{axis}[otSchema, axis lines=middle, width=54mm, height=40mm, xmin=-2, xmax=2, ymin=0, ymax=2,
  xtick=\empty, ytick=\empty, xlabel={$x$}, xlabel style={at={(1,0)},anchor=west}]
  \addplot[otNavy,line width=1.1pt,domain=-1.9:1.9,samples=3] {abs(x)};
  \fill[otNavy] (axis cs:0,0) circle (1.5pt);
\end{axis}
\end{tikzpicture}\\[1mm]
\legende{$x\mapsto|x|$ : continue en $0$,\\ mais non dérivable en $0$}
\end{minipage}

\begin{propriete}[Fonctions continues usuelles]
Les fonctions polynômes, la fonction exponentielle et la valeur absolue sont continues sur $\R$ ; la fonction racine carrée
est continue sur $\intFO{0}{+\infty}$ ; une fonction rationnelle est continue sur chaque intervalle de son ensemble de
définition. Les sommes, produits, quotients et composées de fonctions continues sont continues.
\end{propriete}

\begin{theoreme}[Dérivabilité et continuité]
Toute fonction dérivable sur un intervalle $I$ est continue sur $I$. La réciproque est fausse : $x\mapsto|x|$ est
continue en $0$ mais n'y est pas dérivable.
\end{theoreme}

\begin{savoirfaire}[Prolonger une fonction par continuité]
Si $f$ n'est pas définie en $a$ mais admet une limite finie $\ell$ en $a$, la fonction $g$ définie par $g(x)=f(x)$ pour
$x\neq a$ et $g(a)=\ell$ est continue en $a$ : c'est le \textbf{prolongement par continuité} de $f$ en $a$.
\end{savoirfaire}

\essSection{2}{Le théorème des valeurs intermédiaires}

\begin{theoreme}[Théorème des valeurs intermédiaires]
Soit $f$ une fonction continue sur un intervalle $\intFF{a}{b}$. Pour tout réel $k$ compris entre $f(a)$ et $f(b)$,
l'équation $f(x)=k$ admet \textbf{au moins une} solution dans $\intFF{a}{b}$.
\end{theoreme}

\noindent
\begin{minipage}[t]{0.48\linewidth}\centering
\begin{tikzpicture}
\begin{axis}[otSchema, width=74mm, height=48mm, xmin=0, xmax=6.3, ymin=0, ymax=3.4,
  xtick={0.8,5.2}, xticklabels={$a$,$b$}, ytick={1.6}, yticklabels={$k$},
  xlabel={$x$}, xlabel style={at={(1,0)},anchor=west}]
  \draw[otGreen,dashed,line width=0.6pt] (axis cs:0,1.6) -- (axis cs:6.3,1.6);
  \addplot[otNavy,line width=1.1pt,domain=0.8:5.2,samples=100] {1.6+0.25*(x-3)*((x-3)^2-2.56)};
  \addplot[only marks,mark=*,mark size=1.6pt,otGreen] coordinates {(1.4,1.6) (3,1.6) (4.6,1.6)};
  \draw[otGreen,densely dotted,line width=0.5pt] (axis cs:1.4,1.6) -- (axis cs:1.4,0);
  \draw[otGreen,densely dotted,line width=0.5pt] (axis cs:3,1.6) -- (axis cs:3,0);
  \draw[otGreen,densely dotted,line width=0.5pt] (axis cs:4.6,1.6) -- (axis cs:4.6,0);
  \node[anchor=north,font=\scriptsize,text=otGreenDark] at (axis cs:1.4,0) {$c_{1}$};
  \node[anchor=north,font=\scriptsize,text=otGreenDark] at (axis cs:3,0) {$c_{2}$};
  \node[anchor=north,font=\scriptsize,text=otGreenDark] at (axis cs:4.6,0) {$c_{3}$};
\end{axis}
\end{tikzpicture}\\[0.5mm]
\legende{$f$ continue : au moins une solution de $f(x)=k$.}
\end{minipage}\hfill
\begin{minipage}[t]{0.48\linewidth}\centering
\begin{tikzpicture}
\begin{axis}[otSchema, width=74mm, height=48mm, xmin=0, xmax=6.3, ymin=0, ymax=3.4,
  xtick={0.5,5.8}, xticklabels={$a$,$b$}, ytick={1.7}, yticklabels={$k$},
  xlabel={$x$}, xlabel style={at={(1,0)},anchor=west}]
  \draw[otGreen,dashed,line width=0.6pt] (axis cs:0,1.7) -- (axis cs:6.3,1.7);
  \addplot[otNavy,line width=1.1pt,domain=0.5:5.8,samples=100] {0.3+2.8/(1+exp(-1.3*(x-3.2)))};
  \fill[otGreen] (axis cs:3.2,1.7) circle (1.6pt);
  \draw[otGreen,densely dotted,line width=0.5pt] (axis cs:3.2,1.7) -- (axis cs:3.2,0);
  \node[anchor=north,font=\scriptsize,text=otGreenDark] at (axis cs:3.2,0) {$c$};
\end{axis}
\end{tikzpicture}\\[0.5mm]
\legende{$f$ continue et strictement croissante : une seule solution.}
\end{minipage}

\begin{theoreme}[Cas d'une fonction strictement monotone]
Si $f$ est continue et \textbf{strictement monotone} sur $\intFF{a}{b}$, alors, pour tout réel $k$ compris entre $f(a)$
et $f(b)$, l'équation $f(x)=k$ admet une \textbf{unique} solution dans $\intFF{a}{b}$. Le résultat s'étend à un intervalle
quelconque, en remplaçant $f(a)$ et $f(b)$ par les limites aux bornes.
\end{theoreme}

\begin{savoirfaire}[Montrer qu'une équation $f(x)=k$ a une unique solution]
\begin{enumerate}[label=\otLabel{\arabic*.},leftmargin=1.6em,itemsep=0.6mm,topsep=0.5mm]
  \item Justifier que $f$ est continue (souvent : $f$ est dérivable) sur l'intervalle étudié.
  \item Établir que $f$ est strictement monotone (signe de $f'$), à l'aide du tableau de variations.
  \item Vérifier que $k$ est compris entre les valeurs (ou les limites) aux bornes, puis conclure.
\end{enumerate}
\end{savoirfaire}

\essSection{3}{Encadrer une solution : la dichotomie}

\noindent
\begin{minipage}[c]{0.5\linewidth}
Pour approcher la solution $\alpha$ de $f(x)=0$ sur $\intFF{a}{b}$ (avec $f(a)$ et $f(b)$ de signes contraires), on coupe
l'intervalle en deux à chaque étape et l'on garde la moitié dans laquelle $f$ change de signe. L'amplitude est divisée
par $2$ à chaque étape.
\end{minipage}\hfill
\begin{minipage}[c]{0.45\linewidth}
\centering
\begin{tikzpicture}[x=59mm,y=6mm]
  \foreach \a/\b/\r in {0/1/0,0.5/1/1,0.5/0.75/2,0.625/0.75/3,0.625/0.6875/4}{
    \draw[otRule,line width=0.5pt] (0,-\r) -- (1,-\r);
    \draw[otNavy,line width=2.2pt,line cap=round] (\a,-\r) -- (\b,-\r);
    \node[anchor=east,font=\otLight\fontsize{7}{8}\selectfont,text=otMuted] at (-0.03,-\r) {étape \r};}
  \draw[otGreen,dashed,line width=0.6pt] (0.6823,0.6) -- (0.6823,-4.4);
  \node[anchor=south,font=\scriptsize,text=otGreenDark] at (0.6823,0.6) {$\alpha$};
\end{tikzpicture}
\end{minipage}

\begin{python}
def dichotomie(a, b, p):
    while b - a > 10**(-p):       # amplitude souhaitée : 10^(-p)
        m = (a + b) / 2
        if f(a) * f(m) <= 0:      # changement de signe sur [a ; m]
            b = m
        else:
            a = m
    return a, b
\end{python}

\essSection{4}{Suites et continuité}

\begin{theoreme}[Image d'une suite convergente]
Si la suite $(u_n)$ converge vers $\ell$ et si $f$ est continue en $\ell$, alors la suite $\big(f(u_n)\big)$ converge vers $f(\ell)$.
\end{theoreme}

\begin{theoreme}[Théorème du point fixe]
Soit $(u_n)$ définie par $u_{n+1}=f(u_n)$, avec $f$ continue. Si $(u_n)$ converge vers $\ell$, alors $\ell=f(\ell)$ :
la limite est un \textbf{point fixe} de $f$, c'est-à-dire l'abscisse d'un point d'intersection de $\Cf$ et de la droite $y=x$.
\end{theoreme}

\begin{attention}[Pièges fréquents]
\begin{itemize}[label={\color{otOrange}\textbullet},leftmargin=1.4em,itemsep=0.6mm,topsep=0.5mm]
  \item Le théorème des valeurs intermédiaires donne l'\emph{existence} ; c'est la stricte monotonie qui donne l'\emph{unicité}.
  \item Il faut une fonction continue \emph{sur un intervalle} : $\frac1x$ change de signe sur $\intFF{-1}{1}$ sans jamais s'annuler.
  \item Le point fixe ne s'applique qu'\emph{après} avoir prouvé la convergence ; s'il y a plusieurs points fixes, il faut choisir le bon (encadrement).
  \item Une fonction continue n'est pas forcément dérivable.
\end{itemize}
\end{attention}

%% =====================================================================
%%  EXERCICES
%% =====================================================================
\exercices
\rubrique{Découverte}

\begin{exercice}[1]{Vérifier la continuité}
Étudier la continuité de chacune des fonctions suivantes au point indiqué.
\begin{questions}
  \item $f(x)=\dfrac{x^{2}-4}{x-2}$ en $2$.
  \item $g(x)=x^{2}$ si $x<1$ et $g(x)=2x-1$ si $x\geq1$, en $1$.
  \item $h(x)=\dfrac{\sin x}{x}$ en $0$ (on admet que $\displaystyle\lim_{x\to0}\frac{\sin x}{x}=1$).
  \item $k(x)=\dfrac{x^{2}-1}{x-1}$ si $x\neq1$ et $k(1)=2$, en $1$.
  \item $m(x)=|x-3|$ en $3$.
\end{questions}
\end{exercice}

\rubrique{Entraînement}

\begin{exercice}[2]{Prolongement par continuité}
Chacune des fonctions suivantes n'est pas définie au point indiqué. Montrer qu'elle est prolongeable par continuité
en ce point et préciser la valeur du prolongement.
\begin{tasks}(3)
  \task $\dfrac{x^{3}-8}{x-2}$ en $2$
  \task $\dfrac{\sin(2x)}{x}$ en $0$
  \task $\dfrac{e^{2x}-1}{x}$ en $0$
  \task $\dfrac{e^{x}-1}{x}$ en $0$
  \task $\dfrac{\sqrt{x+1}-1}{x}$ en $0$
\end{tasks}
\end{exercice}

\begin{exercice}[2]{Différents types de discontinuités}
Pour chacune des fonctions suivantes, étudier les limites à gauche et à droite au point indiqué, puis dire si la fonction
est continue, prolongeable par continuité, ou ni l'un ni l'autre.
\begin{questions}
  \item $a(x)=\dfrac{1}{x-1}$ en $1$.
  \item $g(x)=x+1$ si $x<0$ et $g(x)=x-1$ si $x\geq0$, en $0$.
  \item $h(x)=\dfrac{1}{x^{2}}$ en $0$.
  \item $k(x)=\sin\left(\dfrac1x\right)$ si $x\neq0$ et $k(0)=0$, en $0$. On pourra étudier $k(x_n)$ et $k(y_n)$ avec
    $x_n=\dfrac{1}{2n\pi+\frac{\pi}{2}}$ et $y_n=\dfrac{1}{2n\pi-\frac{\pi}{2}}$.
  \item $m(x)=x^{2}$ si $x<1$ et $m(x)=x+1$ si $x\geq1$, en $1$.
\end{questions}
\end{exercice}

\begin{exercice}[2]{Théorème des valeurs intermédiaires}
\begin{questions}
  \item Montrer que l'équation $x^{3}-3x+1=0$ admet au moins une solution dans $\intFF{0}{1}$ et au moins une dans $\intFF{1}{2}$.
  \item L'équation $e^{x}=2x$ admet-elle une solution dans $\intFF{0}{1}$ ? On pourra utiliser l'inégalité $e^{x}\geq1+x$.
  \item Montrer que l'équation $e^{-x}=x$ admet au moins une solution dans $\intFF{0}{1}$.
  \item Montrer que l'équation $\cos x=x$ admet au moins une solution dans $\intFF{0}{\frac{\pi}{2}}$.
  \item Peut-on utiliser le théorème des valeurs intermédiaires pour résoudre $x^{4}-2x^{2}+1=0$ ? Résoudre cette équation.
\end{questions}
\end{exercice}

\rubrique{Problèmes}

\begin{exercice}[3]{Continuité, paramètres et nombre de solutions}
\begin{questions}
  \item Déterminer le réel $a$ pour que la fonction définie par $f(x)=\dfrac{x^{2}-4}{x-2}$ si $x\neq2$ et $f(2)=a$ soit continue en $2$.
  \item Montrer que l'équation $x^{5}-5x+1=0$ admet exactement trois solutions réelles.
  \item On considère la fonction $g$ définie par $g(x)=ax+b$ si $x<1$ et $g(x)=x^{2}$ si $x\geq1$.
    \begin{questions}
      \item À quelle condition sur $a$ et $b$ la fonction $g$ est-elle continue sur $\R$ ?
      \item Déterminer $a$ et $b$ pour que $g$ soit de plus dérivable en $1$.
    \end{questions}
  \item Montrer que l'équation $e^{x}=x+2$ admet exactement deux solutions réelles.
  \item Déterminer le réel $k$ pour que la fonction définie par $h(x)=\dfrac{\sin x}{x}$ si $x\neq0$ et $h(0)=k$ soit continue en $0$.
\end{questions}
\end{exercice}

\begin{exercice}[3]{Une fonction auxiliaire}
\partieExo{A}{Étude d'une fonction auxiliaire}
On considère la fonction $g$ définie sur $\R$ par $g(x)=x^{3}+x-1$.
\begin{questions}
  \item Calculer les limites de $g$ en $-\infty$ et en $+\infty$, puis étudier ses variations.
  \item Montrer que l'équation $g(x)=0$ admet une unique solution $\alpha$ sur $\R$, et que $\alpha\in\intFF{0}{1}$.
  \item On considère la fonction Python ci-dessous. Que renvoie \lstinline[style=otPython]!dichotomie(2)! ? En déduire un
    encadrement de $\alpha$ d'amplitude $10^{-2}$.
\begin{python}
def dichotomie(p):
    a, b = 0, 1
    while b - a > 10**(-p):
        m = (a + b) / 2
        if m**3 + m - 1 < 0:
            a = m
        else:
            b = m
    return a, b
\end{python}
  \item En déduire le signe de $g(x)$ selon les valeurs de $x$.
\end{questions}
\partieExo{B}{Étude d'une fonction}
On considère la fonction $f$ définie sur $\R$ par $f(x)=\dfrac{x^{4}}{4}+\dfrac{x^{2}}{2}-x$.
\begin{questions}[resume]
  \item Calculer $f'(x)$ et en déduire les variations de $f$.
  \item Montrer que $f(\alpha)=\dfrac{\alpha(\alpha-3)}{4}$. On pourra utiliser l'égalité $\alpha^{3}=1-\alpha$.
  \item En déduire un encadrement du minimum de $f$.
\end{questions}
\end{exercice}

\begin{exercice}[3]{Suite récurrente et point fixe}
On considère la fonction $f$ définie sur $\R$ par $f(x)=\dfrac{x^{2}+4}{5}$ et la suite $(u_n)$ définie par $u_0=2$ et,
pour tout entier naturel $n$, $u_{n+1}=f(u_n)$.

\noindent
\begin{minipage}[c]{0.56\linewidth}
\begin{questions}
  \item Étudier les variations de $f$ sur $\intFO{0}{+\infty}$.
  \item Sur la figure ci-contre, construire $u_1$, $u_2$ et $u_3$. Conjecturer le comportement de la suite.
  \item Démontrer par récurrence que, pour tout entier naturel $n$, $1\leq u_{n+1}\leq u_n\leq2$.
  \item En déduire que $(u_n)$ converge vers un réel $\ell\in\intFF{1}{2}$.
\end{questions}
\end{minipage}\hfill
\begin{minipage}[c]{0.4\linewidth}
\centering
\begin{tikzpicture}
\begin{axis}[otGraphe, width=58mm, xmin=0, xmax=2.3, ymin=0, ymax=2.3,
  xtick={1,2}, ytick={1,2}, minor tick num=4, grid=both, minor grid style={otRule!45,line width=0.2pt}]
  \addplot[otRule!80!otMuted,line width=0.6pt,domain=0:2.25] {x};
  \addplot[otNavy,line width=1.1pt,domain=0:2.25,samples=60] {(x^2+4)/5};
  \addplot[only marks,mark=*,mark size=1.4pt,otGreen] coordinates {(2,0)};
  \node[anchor=north,font=\scriptsize] at (axis cs:2,-0.08) {$u_0$};
  \node[anchor=south east,font=\scriptsize,text=otNavy] at (axis cs:0.6,0.95) {$\Cf$};
\end{axis}
\end{tikzpicture}
\end{minipage}
\begin{questions}[start=5]
  \item Justifier que $\ell=f(\ell)$, puis déterminer $\ell$.
  \item À l'aide de la calculatrice, déterminer le plus petit entier $n$ tel que $u_n-1<10^{-3}$.
  \item Que se passe-t-il si l'on choisit $u_0=4$ ? Et si l'on choisit $u_0=5$ ?
\end{questions}
\end{exercice}

\begin{exercice}[3]{Trois démonstrations d'existence}
\partieExo{A}{Un point fixe}
Soit $f$ une fonction continue sur $\intFF{0}{1}$, à valeurs dans $\intFF{0}{1}$.
\begin{questions}
  \item On pose $g(x)=f(x)-x$. Déterminer le signe de $g(0)$ et de $g(1)$.
  \item En déduire qu'il existe $c\in\intFF{0}{1}$ tel que $f(c)=c$. Interpréter graphiquement.
\end{questions}
\partieExo{B}{Le randonneur}
Un randonneur quitte un refuge à 8~h et atteint un sommet à 16~h. Le lendemain, il redescend par le même chemin, de 8~h à
16~h. On note $d_1(t)$ et $d_2(t)$ la distance (en km) qui le sépare du refuge, le premier et le second jour, à l'heure $t$ ;
ces fonctions sont continues sur $\intFF{8}{16}$, et la longueur du chemin est $L$.
\begin{questions}[resume]
  \item Montrer qu'il existe une heure à laquelle le randonneur se trouve au même endroit les deux jours.
\end{questions}
\partieExo{C}{La coureuse}
Une coureuse parcourt $10$~km en $50$~minutes, à une allure variable. On note $d(t)$ la distance parcourue (en km) au bout de
$t$ minutes ; la fonction $d$ est continue sur $\intFF{0}{50}$, avec $d(0)=0$ et $d(50)=10$. Pour $t\in\intFF{0}{25}$, on pose
$h(t)=d(t+25)-d(t)-5$.
\begin{questions}[resume]
  \item Calculer $h(0)+h(25)$.
  \item En déduire qu'il existe une période de $25$ minutes pendant laquelle la coureuse parcourt exactement $5$~km.
\end{questions}
\end{exercice}
